\chapter{部署：推理框架与生成服务}
\label{aigc:deployment}

\section{按工作负载选择框架}

AIGC 服务通常包含两类计算：脚本、提示改写和素材理解等语言/视觉语言任务，
以及图像、视频、音频的生成任务。前者可按模型兼容性选择 vLLM 或 SGLang，
利用连续批处理、前缀复用等能力提高服务效率；后者应选择支持对应生成模型的执行管线。

当前 SGLang 已提供 SGLang Diffusion，vLLM 生态也有 vLLM-Omni，
覆盖部分扩散及多模态生成任务。因此不能把这两个生态只理解为文本服务，
也不能因为框架支持 DiT，就认为任意图像或视频权重都能直接接入。
实际选型需同时核对模型、codec、调度器、精度、硬件与所用发行版本。

TensorRT 可作为特定硬件上的编译与执行优化工具，Triton Inference Server
可用于组织模型服务；它们与请求调度、生成流水线不是同一层次。
首版优先跑通官方支持的模型与推理路径，再根据测得的瓶颈决定是否增加这些组件。

\section{先建立质量基线，再优化性能}

对同一模型固定分辨率、时长、参考数量、采样器、步数和引导强度，
分别测量排队、条件编码、去噪、解码与输出耗时，同时记录峰值显存、
吞吐和端到端延迟的中位数及95分位数。语言模型还应区分首 token 延迟与后续生成速度。
不能只报一个最快请求，或把文本服务的吞吐直接当作视频生成效率。

优化通常从按形状分桶、批处理、静态条件缓存和阶段卸载开始，
再考虑量化、编译、稀疏或少步模型。每次只改变少量因素，并检查细节、
条件遵循和稳定性。普通双向扩散的目标状态随去噪步变化，
不能直接套用自回归语言模型的跨步 KV 缓存假设。

\section{服务接口应保留哪些信息}

每次请求保存模型与配置版本、随机种子、素材版本、生成参数和输出记录，
提供队列、取消、超时与失败重试；重试产生的费用也计入任务。
涉及局部编辑时，明确保留区域、重做范围及音轨处理规则，
才能在生成失败后定向修正。部署的目标是稳定地产出可验收结果，
因此还需统计每条验收结果的总成本，而不只统计单次调用价格。

\par\medskip
{\small 本节框架能力核查于2026年11月24日，依据官方
\href{https://github.com/vllm-project/vllm}{vLLM}、
\href{https://github.com/sgl-project/sglang}{SGLang} 与
\href{https://github.com/vllm-project/vllm-omni}{vLLM-Omni} 的项目说明。
具体模型与硬件支持以所选版本为准；本文未作性能实测。\par}

\chapter{应用与商业化：设计、短视频与短剧}
\label{aigc:applications}

\section{以生产任务组织能力}

商业应用应从客户需要交付什么出发，而不是按模型品牌列目录。
可先围绕以下三类场景组织资产、工具和验收规则。

\begin{itemize}
\item \textbf{设计。}用于商品视觉、广告创意、海报与场景草图，
重点是资产检索、多方案比较、局部改稿和尺寸适配。
品牌字样、商品结构、价格和最终版式由明确素材与排版环节控制，
不能只按画面是否好看验收。
\item \textbf{短视频。}将脚本、商品或人物素材转成镜头，再完成剪辑、配音、
字幕与混流。重点是前几秒表达、镜头可用性、声画关系与版本复用；
先做短镜头与定向返工，再扩展时长。
\item \textbf{短剧与叙事内容。}红果等平台的微短剧观看与发行场景，
可作为确定受众、竖屏规格和叙事节奏的应用参照。
生产端需要角色与声音资产、分集脚本、分镜、服装道具状态和镜头连续性。
平台示例不等于其已接收任意 AI 成片；具体合作、内容要求与分账条件应另行核实。
\end{itemize}

\section{最小可交付流程}

一个可用流程是：需求与素材确认，形成脚本或设计方案，按镜头/区域生成候选，
人工选择与局部修改，完成确定性排版或后期，最后按交付规格验收。
客户确认的商品、角色、文案和已通过镜头应保留版本，
后续修改尽量只重做失败部分，避免每次整条推倒重来。

设计项目记录返修次数与人工时间；短视频记录可用镜头比例、剪辑时间与验收率；
短剧增加跨镜头身份、对白、道具状态和剧情连贯性的检查。
上线后的点击、完播、转化或付费表现还受选题、投放与分发影响，
应与生成质量分别记录。生成速度变快不自动等于业务收益增长。

\section{先证明交付，再决定是否自建模型}

首个项目宜选择范围明确、素材可复用的任务，例如商品设计改稿或短视频镜头生产。
先用现有模型与外部工具跑通小规模交付，完整核算生成、失败候选、
人工审片、修改、后期和服务费用。
只有当同类需求持续出现、失败集中且现有接口确实不足时，
再投入专用数据、适配模块或模型训练。是否自建，取决于可替代的交付量与净节省，
而不只取决于单次推理是否便宜。
